% ============================ CHAPITRE 3 =============================
\chapter{Réalisation, implémentation et validation}
\label{chap:realisation}

\section{Introduction}

Ce dernier chapitre décrit la mise en œuvre effective de l'architecture conçue au
chapitre~\ref{chap:conception}. Il présente l'environnement de développement (dont les
contraintes ont pesé sur la réalisation autant que le cahier des charges lui-même),
détaille l'implémentation des composants clés et l'intégration des services tiers, expose
les mécanismes de sécurité applicative, puis la démarche de validation : évaluation
quantitative RAGAS, scénario de démonstration, résultats obtenus, difficultés rencontrées
et limites de la solution.

\section{Environnement de développement}
\label{sec:env-dev}

\subsection{Poste de travail et contraintes}

Le développement s'est déroulé sur un poste professionnel ALTEN sous Windows 11,
\emph{verrouillé} : pas de droits d'administration, pas d'installation d'applications ou
de services système (Docker, serveurs de bases de données, Tesseract...), et un agent
\textbf{Zscaler Client Connector} interceptant l'intégralité du trafic HTTPS sortant par
inspection TLS, quel que soit le réseau utilisé, l'agent étant applicatif et non
réseau.

\begin{alertbox}[Conséquences concrètes des contraintes]
\begin{itemize}
  \item Toute dépendance devait être installable \emph{sans droits administrateur} :
        environnement virtuel Python (\texttt{venv}), paquets \texttt{pip} et \texttt{npm}
        locaux au projet ;
  \item Aucun service ne pouvait être installé : d'où SQLite, ChromaDB embarqué et
        FakeRedis en processus ;
  \item L'inspection TLS de Zscaler re-signe le trafic HTTPS avec son propre certificat
        racine, reconnu par Windows mais pas par la liste \texttt{certifi} embarquée de
        Python : toutes les requêtes HTTPS Python échouaient en
        \texttt{CERTIFICATE\_VERIFY\_FAILED} alors que le navigateur fonctionnait ;
  \item Certains domaines étaient bloqués : le CDN de stockage de Hugging Face
        (téléchargement de modèles) et son API d'inférence, ce qui a imposé des
        stratégies de contournement décrites plus loin.
\end{itemize}
\end{alertbox}

\subsection{Outillage}

L'environnement s'appuie sur : Python 3.12 en environnement virtuel, Node.js/npm pour le
frontend (Vite assurant serveur de développement et rechargement à chaud), Git/GitHub pour
le versionnement, Visual Studio Code assisté d'outils d'IA de développement, et des
notebooks Jupyter pour les expérimentations préalables (validation du parsing Docling sur
\emph{Attention Is All You Need}, comparaison de VLM pour le captioning, campagne
d'évaluation RAGAS). L'API est documentée automatiquement par le Swagger UI de FastAPI
(\texttt{/docs}), qui a servi de banc de test permanent.

\section{Implémentation des composants principaux}
\label{sec:implementation}

\subsection{Analyse documentaire et découpage (\texttt{docling\_parser})}

Le convertisseur Docling est construit une seule fois (mise en cache par
\texttt{lru\_cache}) avec des options dédiées : génération des images de figures activée,
OCR désactivé (hors périmètre de cette itération), backend PDF \texttt{pypdfium2} et,
lorsque la variable \texttt{DOCLING\_ARTIFACTS\_PATH} est renseignée, chargement des
modèles de layout et de reconstruction de tableaux depuis un dossier local plutôt que
depuis Hugging Face : le mécanisme qui a permis de fonctionner malgré le blocage du CDN.

Le découpage utilise le \emph{HybridChunker} de \texttt{docling-core}, paramétré par le
tokenizer du modèle d'embeddings (512 tokens max) ; chaque chunk est \emph{contextualisé}
(préfixé des titres de sections englobantes) et rattaché à son numéro de page pour la
citation. Un découpage de repli à fenêtre glissante (1500 caractères, chevauchement de
200) prend le relais si le chunker sémantique échoue sur un document atypique.

\subsection{Description des figures par VLM (\texttt{hf\_client})}

Chaque figure extraite est envoyée au VLM \textbf{Gemma-3 27B} (servi par DeepInfra via
l'API unifiée Hugging Face Inference Providers) sous forme d'URL de données base64, avec
un prompt d'ancrage soigneusement construit : nom du document et légende d'origine de la
figure fournis comme contexte, consigne de décrire type, composants, relations, textes et
valeurs visibles, interdiction des préambules (« Voici une description... ») et du
\emph{hedging} sur les faits déjà donnés. La description finale indexée est la
concaténation \emph{légende d'origine + description VLM} : la légende (« Figure 1: The
Transformer architecture ») donne une prise à la recherche lexicale exacte, la
description VLM la profondeur sémantique.

Les appels aux API d'inférence sont protégés par une politique de nouvelle tentative
sélective (bibliothèque \texttt{tenacity}) : trois essais à attente exponentielle
aléatoire sur les erreurs 5xx transitoires (démarrages à froid des fournisseurs
serverless), aucune tentative sur les 4xx (erreur d'authentification ou de charge utile,
qu'un nouvel essai ne réparera pas).

\subsection{Embeddings multi-backends}

Le module d'embeddings illustre l'adaptation aux contraintes réseau : trois backends
commutables par la seule variable \texttt{EMBEDDING\_BACKEND}, derrière une interface
unique \texttt{embed\_texts()} :

\begin{itemize}
  \item \texttt{hf\_inference} : \texttt{multilingual-e5-large} via l'API Hugging Face,
        le choix initial, devenu inutilisable lorsque le domaine a été bloqué ;
  \item \texttt{local} : les mêmes poids e5-large exécutés sur CPU via
        \texttt{transformers}, chargés depuis un dossier local en
        \texttt{local\_files\_only} (jamais de réseau), avec \emph{mean pooling} masqué et
        normalisation L2 reproduisant exactement la recette du serveur d'inférence HF,
        garantissant des vecteurs compatibles avec l'index existant ;
  \item \texttt{mistral} : \texttt{mistral-embed} via l'API Mistral, accessible à travers
        Zscaler, le backend actif. Les vecteurs (1024 dim.) sont normalisés L2 pour que
        la distance de ChromaDB ordonne comme une similarité cosinus, et envoyés par lots
        de 32 (une requête pour plusieurs dizaines de textes, compatible avec la limite de
        débit du palier gratuit).
\end{itemize}

\begin{alertbox}[Un piège d'ingénierie : deux espaces vectoriels de même dimension]
\texttt{mistral-embed} et e5-large produisent tous deux des vecteurs de 1024 dimensions
mais dans des espaces \emph{incompatibles}. ChromaDB ne signale donc aucune erreur si l'on
change de backend sans réindexer : il renvoie silencieusement des similarités absurdes.
La bascule vers Mistral a ainsi nécessité la purge de la collection et la réingestion
complète du corpus, et le risque est explicitement documenté dans la configuration.
\end{alertbox}

\subsection{Recherche hybride (\texttt{vector\_store})}

La recherche dense interroge ChromaDB avec filtres de métadonnées optionnels (liste de
documents, type de chunk). La recherche lexicale s'appuie sur un index BM25 en mémoire
reconstruit paresseusement depuis la collection ChromaDB (source de vérité unique)
lorsqu'il est marqué obsolète par une ingestion ou une suppression ; ce choix évite un
second index persistant à maintenir synchrone, pour un coût de reconstruction négligeable
à l'échelle du corpus. La fusion RRF additionne, pour chaque chunk, les
\(1/(60 + \mathrm{rang})\) de ses positions dans les deux classements, sur un
sur-échantillonnage de \(3 \times k\) candidats par voie ; les scores fusionnés sont
normalisés par le maximum pour l'affichage.

\subsection{Routage automatique du périmètre (\texttt{auto-scope})}

Lorsque l'utilisateur laisse le périmètre en automatique, un appel LLM léger
(température 0) reçoit le catalogue \emph{groupes → documents indexés} et la question, et
répond en JSON strict par la liste des groupes plausiblement pertinents. La conception est
volontairement \emph{conservatrice} : toute défaillance (réponse non parsable, groupe
inconnu, exception réseau) ou toute réponse équivalente à « tous les groupes » désactive
le filtre : le routage ne peut jamais rétrécir le corpus par accident, seulement l'élaguer
quand le modèle est confiant. Les groupes retenus sont affichés à l'utilisateur avec la
réponse (« recherché dans : ... »), pour la transparence du périmètre.

\subsection{Génération contrainte et restitution multimodale}

Le prompt système de génération impose : réponse fondée \emph{uniquement} sur les sources
numérotées, citations \texttt{[Source N]} en ligne, aveu explicite lorsque les sources ne
contiennent pas la réponse. Il traite aussi une subtilité multimodale : le modèle est
informé que certaines sources sont des figures (dont il ne voit que la description) et que
l'application affiche automatiquement l'image citée sous la réponse ; sans cette
instruction, le LLM s'excusait de « ne pas pouvoir montrer d'images ». Côté API, chaque
source est restituée typée (\texttt{text} avec extrait, ou \texttt{image} avec URL et
description), et le frontend affiche les cartes de sources et les figures citées.

\subsection{Interface utilisateur}

L'interface React comprend trois espaces : l'accueil, l'espace \emph{Documents} (dropzone
de téléversement multiple, cartes de documents avec état d'ingestion en direct (étape
courante affichée pendant le traitement), vue détail avec rendu Markdown et galerie de
figures, gestion des groupes) et l'espace \emph{Chat} (historique des conversations,
sélecteur de périmètre documentaire, rendu Markdown des réponses avec
\texttt{react-markdown}, cartes de sources cliquables, visionneuse de figures). Le thème
clair/sombre est persistant et l'ensemble des appels passe par un client API central qui
joint le jeton de session Clerk.

\section{Intégration des services tiers}
\label{sec:integrations}

\begin{table}[!htb]
\centering\small
\caption{Services externes intégrés et leur rôle.}
\label{tab:services}
\begin{tabularx}{\textwidth}{@{}p{3.4cm}p{4.2cm}X@{}}
\toprule
\textbf{Service} & \textbf{Rôle} & \textbf{Modalités d'intégration} \\
\midrule
Mistral AI & Embeddings + génération + routage & SDK \texttt{mistralai} ; batch d'embeddings normalisés L2 ; température 0 pour le routage JSON \\
HF Inference Providers (DeepInfra) & Captioning VLM Gemma-3 27B & \texttt{huggingface\_hub.InferenceClient} ; images en data-URL base64 ; retry sélectif 5xx \\
Clerk & Authentification & SDK React côté client ; côté backend, vérification RS256 pure via JWKS (aucun appel API Clerk par requête) \\
Groq & Juge LLM d'évaluation (Llama 3.3 70B) & Utilisé uniquement dans le notebook RAGAS, indépendant du système évalué \\
Redis / OCI Cache & Cache de lecture des conversations & Client \texttt{redis} ou \texttt{fakeredis} selon \texttt{REDIS\_URL} ; fail-open intégral \\
\bottomrule
\end{tabularx}
\end{table}

Un détail d'intégration élégant côté Clerk : la clé publique (\texttt{pk\_...}) encode en
base64 le domaine du Frontend API, qui est aussi l'émetteur des jetons et l'hôte du JWKS ;
le backend décode ce domaine au démarrage et n'a besoin d'aucune configuration
supplémentaire ni d'appel de découverte.

\section{Détails techniques spécifiques à l'environnement contraint}
\label{sec:details-contraintes}

\subsection{Traversée du proxy Zscaler : \texttt{truststore}}

Le problème des erreurs \texttt{CERTIFICATE\_VERIFY\_FAILED} (section~\ref{sec:env-dev})
a été résolu proprement par le paquet \texttt{truststore} : un appel unique
\texttt{truststore.inject\_into\_ssl()} à l'import du paquet applicatif, donc avant tout
appel réseau, fait utiliser à Python le magasin de certificats de Windows (qui contient
les racines ALTEN/Zscaler) au lieu de la liste \texttt{certifi} embarquée. Aucune
désactivation de vérification TLS, aucune variable d'environnement fragile : la solution
est sûre et transparente pour tout le code aval.

\subsection{Modèles Docling hors ligne}

Le CDN de Hugging Face étant bloqué, les modèles de Docling (détection de layout,
TableFormer) ont été téléchargés sur une machine non restreinte
(\texttt{docling-tools models download}), archivés, publiés en artefact de release GitHub
(domaine autorisé), puis dépointés sur le poste via \texttt{DOCLING\_ARTIFACTS\_PATH}. Deux
guides de mise en place (poste d'entreprise / poste personnel) documentent la procédure
pour les repreneurs du projet.

\subsection{Cache de conversations fail-open}

Le cache JSON des lectures de conversations est conçu pour ne \emph{jamais} pouvoir nuire :
toute erreur (Redis injoignable, valeur corrompue, exception du client) est interprétée
comme un cache absent et l'application continue sur SQLite. Le backend est choisi par
\texttt{REDIS\_URL} : \texttt{memory://} (FakeRedis en processus, poste verrouillé),
URL Redis réelle (OCI Cache en production), ou vide (désactivé).

\section{Sécurité applicative}
\label{sec:securite}

Sans viser la gestion fine des droits (hors périmètre), le système applique les
fondamentaux :

\begin{itemize}
  \item \textbf{Authentification forte déléguée} : sessions gérées par Clerk ; le backend
        vérifie chaque jeton par signature RS256 contre le JWKS de l'émetteur (clés
        publiques mises en cache), contrôle l'émetteur et l'expiration, et rejette en 401
        tout jeton absent ou invalide ;
  \item \textbf{Cloisonnement des données} : toute lecture/écriture de conversation vérifie
        que la conversation appartient à l'utilisateur du jeton (réponse 404 sinon, sans
        révéler l'existence de la ressource) ;
  \item \textbf{Validation des entrées} : schémas Pydantic sur toutes les routes, liste
        blanche d'extensions de fichiers au téléversement, validation des noms de groupes,
        identifiants générés côté serveur (UUID) ;
  \item \textbf{CORS restrictif} : origines autorisées explicites par configuration ;
  \item \textbf{Secrets hors dépôt} : clés d'API dans \texttt{.env} non versionné, avec
        \texttt{.env.example} documenté ;
  \item \textbf{TLS vérifié de bout en bout} : \texttt{truststore} maintient la
        vérification des certificats à travers le proxy d'inspection (plutôt que de la
        désactiver, réflexe dangereux courant dans ce type d'environnement).
\end{itemize}

\section{Tests et validation : évaluation RAGAS}
\label{sec:evaluation}

\subsection{Méthodologie}

La qualité du système a été évaluée avec \textbf{RAGAS 0.4.3} \cite{es2023ragas} sur un
corpus de trois documents indexés (\emph{BERT}, \emph{Attention Is All You Need},
\emph{Circuit Diagram Symbols}) et six questions de référence rédigées manuellement avec
leurs réponses attendues, couvrant les deux familles documentaires (articles de recherche
et manuel technique riche en figures).

Deux précautions méthodologiques structurent le protocole :

\begin{itemize}
  \item \textbf{Juge indépendant} : les métriques à base de LLM sont calculées par
        \texttt{llama-3.3-70b-versatile} (servi par Groq), un modèle \emph{étranger} au
        système évalué (qui génère avec Mistral), afin d'éviter le biais
        d'auto-évaluation ;
  \item \textbf{Espace vectoriel cohérent} : la métrique \emph{Answer Relevancy}, qui
        compare des embeddings, utilise \texttt{mistral-embed}, les embeddings mêmes du
        système, pour mesurer la pertinence dans l'espace où le système opère.
\end{itemize}

Quatre métriques sont retenues : \emph{Faithfulness} (la réponse est-elle fidèle aux
passages récupérés, sans hallucination ?), \emph{Answer Relevancy} (répond-elle à la
question ?), \emph{Context Precision} (les passages récupérés sont-ils pertinents ?) et
\emph{Context Recall} (l'information nécessaire a-t-elle été récupérée ?) ; les deux
premières évaluant la génération, les deux dernières le retrieval.

\subsection{Résultats obtenus}

\begin{table}[!htb]
\centering\small
\caption{Résultats de l'évaluation RAGAS (juge : Llama 3.3 70B via Groq).}
\label{tab:ragas}
\begin{tabularx}{\textwidth}{@{}Xcccc@{}}
\toprule
\textbf{Question} & \textbf{Faithf.} & \textbf{Ans. Rel.} & \textbf{Ctx. Prec.} & \textbf{Ctx. Rec.} \\
\midrule
Que signifie l'acronyme BERT ? & 1{,}000 & 0{,}890 & 0{,}756 & 1{,}000 \\
Sur quelles deux tâches BERT est-il pré-entraîné ? & 1{,}000 & 0{,}896 & 0{,}700 & 1{,}000 \\
Quel pourcentage de tokens est masqué lors du MLM ? & 1{,}000 & 0{,}959 & 0{,}700 & 1{,}000 \\
Quel token spécial est ajouté en début de séquence ? & 1{,}000 & 0{,}987 & 1{,}000 & 1{,}000 \\
À quoi sert un schéma de circuit électrique ? & 1{,}000 & 0{,}926 & 0{,}888 & 1{,}000 \\
Que représente le symbole de mise à la terre ? & 0{,}000 & 0{,}926 & 1{,}000 & 1{,}000 \\
\midrule
\textbf{Moyenne} & \textbf{0{,}833} & \textbf{0{,}931} & \textbf{0{,}841} & \textbf{1{,}000} \\
\bottomrule
\end{tabularx}
\end{table}

\subsection{Analyse des résultats}

\begin{itemize}
  \item \textbf{Context Recall = 1{,}00 partout} : l'information nécessaire est
        systématiquement récupérée : la validation la plus directe de l'indexation
        multimodale unifiée et de la recherche hybride dense + BM25 ;
  \item \textbf{Answer Relevancy moyenne de 0{,}93} : les réponses adressent bien les
        questions posées ;
  \item \textbf{Faithfulness parfaite sur 5 questions sur 6} ; le zéro de la question sur
        le symbole de mise à la terre est instructif : la réponse décrivait
        l'\emph{apparence} du symbole (fidèle à la description VLM récupérée) plutôt que sa
        \emph{signification électrique} attendue ; le contexte récupéré, issu d'une
        description d'image, ne contenait pas la définition sémantique. L'axe
        d'amélioration identifié porte sur le prompt de captioning et de génération
        (distinguer « décrire » d'« expliquer la signification ») ;
  \item \textbf{Context Precision de 0{,}70--0{,}76 sur les questions BERT} : le retrieval
        remonte parfois des chunks d'\emph{Attention Is All You Need} non pertinents pour
        BERT, les deux articles partagent un vocabulaire très proche. Pistes : re-ranking
        ou filtrage par document.
\end{itemize}

\rapportfig[0.85]{ragas_radar.png}{Synthèse visuelle des quatre métriques RAGAS moyennes
du système.}{fig:ragas-radar}

\section{Scénario de démonstration}
\label{sec:demo}

Le scénario de démonstration de fin de stage déroule le cycle complet :

\begin{enumerate}
  \item Connexion (Clerk), arrivée sur l'espace documentaire vide ;
  \item Téléversement en rafale de plusieurs PDF (article BERT + manuel de schémas
        électriques) ; observation du traitement séquentiel et des étapes d'ingestion en
        direct (« Describing figure 3 of 5... ») ;
  \item Ouverture du détail d'un document : Markdown structuré (tableaux reconstruits),
        galerie des figures avec descriptions générées par le VLM ;
  \item Question textuelle (« Sur quelles tâches BERT est-il pré-entraîné ? ») : réponse
        citée \texttt{[Source 1]}\texttt{[Source 2]} avec extraits et pages ;
  \item Question visuelle (« Montre-moi l'architecture du Transformer ») : la réponse cite
        une source de type figure et l'image s'affiche sous la réponse ;
  \item Question avec périmètre automatique : l'interface indique les groupes retenus par
        le routeur LLM ;
  \item Question hors corpus : le système déclare explicitement ne pas trouver la réponse
        dans les sources : démonstration du garde-fou anti-hallucination.
\end{enumerate}

\section{Résultats obtenus}
\label{sec:resultats}

Au terme des huit semaines, l'ensemble des objectifs du cahier des charges est atteint :

\begin{itemize}
  \item une chaîne d'ingestion multimodale complète et robuste (PDF/Word/Excel,
        figures décrites et indexées, file sérialisée, suivi d'état, gestion d'erreurs) ;
  \item un moteur de restitution hybride (dense + BM25 + RRF) avec routage automatique du
        périmètre et génération contrainte citée ;
  \item une API REST documentée (OpenAPI) d'une quinzaine de routes couvrant documents,
        groupes, conversations et chat ;
  \item une interface web complète et soignée (authentification, gestion du corpus, chat
        multi-conversations, affichage des sources et figures, thème clair/sombre) ;
  \item une évaluation quantitative RAGAS aux résultats solides (recall 1{,}00 ;
        relevancy 0{,}93 ; precision 0{,}84 ; faithfulness 0{,}83) avec analyse des cas
        d'échec ;
  \item une base prête pour le déploiement OCI (stockage embarqué, cache commutable,
        configuration par environnement).
\end{itemize}

\section{Difficultés rencontrées}
\label{sec:difficultes}

\begin{description}
  \item[Blocages réseau successifs.] La difficulté dominante du stage. L'inspection TLS de
    Zscaler a d'abord cassé toutes les requêtes HTTPS Python (résolu par
    \texttt{truststore}) ; puis le blocage du CDN Hugging Face a empêché le téléchargement
    des modèles Docling (résolu par le circuit hors ligne via release GitHub) ; enfin le
    blocage de l'API d'inférence HF a invalidé le backend d'embeddings initial, imposant
    la bascule vers \texttt{mistral-embed}, et donc la purge et la réingestion complète de
    l'index (espaces vectoriels incompatibles à dimension égale) ;
  \item[Robustesse du parsing PDF.] Le backend Docling par défaut s'effondrait en
    \texttt{std::bad\_alloc} sur des PDF à images volumineuses ; le passage au backend
    \texttt{pypdfium2} a réglé le problème au prix d'une fidélité de layout légèrement
    moindre ;
  \item[Contention d'écriture sous téléversements en rafale.] Les ingestions concurrentes
    (une par requête via \texttt{BackgroundTasks}) saturaient le poste et faisaient
    échouer des documents par contention SQLite/ChromaDB, résolu par la file à worker
    unique ;
  \item[Comportement du LLM face aux sources visuelles.] Sans instruction dédiée, le
    modèle s'excusait de ne pas pouvoir afficher d'images (alors que l'application le fait
    pour lui) ; et les descriptions VLM initiales, trop « visuelles », ont motivé le
    prompt d'ancrage contextuel et le préfixage des légendes d'origine ;
  \item[Calendrier contraint.] Huit semaines pour l'ensemble de la chaîne : les choix de
    sobriété (pas de LangChain, stockage embarqué, services managés) ont été autant des
    choix de qualité que des conditions de faisabilité.
\end{description}

\section{Limites de la solution}
\label{sec:limites}

\begin{itemize}
  \item \textbf{Passage à l'échelle du retrieval} : l'index BM25 est reconstruit en
        mémoire depuis ChromaDB, approprié à un corpus de démonstration, à revoir
        (index persistant, base dédiée) pour des milliers de documents ;
  \item \textbf{Pas d'OCR} : les documents scannés (images de texte) ne sont pas couverts
        par cette itération ; l'architecture le permet (option Docling), la contrainte
        d'installation de Tesseract ou le coût d'un service cloud l'ont différé ;
  \item \textbf{Dépendance aux API externes} : embeddings, captioning et génération
        reposent sur des services tiers : disponibilité, latence et limites de débit
        s'imposent au système ;
  \item \textbf{Évaluation à petite échelle} : six questions de référence suffisent à
        diagnostiquer, pas à qualifier statistiquement ; un jeu d'évaluation étendu et
        automatisé serait la suite naturelle ;
  \item \textbf{Corpus partagé entre utilisateurs} : conformément au périmètre, seules les
        conversations sont cloisonnées par utilisateur ; le corpus documentaire est commun
        (pas de gestion fine des droits).
\end{itemize}

\section{Conclusion}

Ce chapitre a montré la traduction de l'architecture en un système opérationnel, validé
quantitativement, et façonné en profondeur par son environnement d'exécution : chaque
contrainte du poste verrouillé a trouvé une réponse d'ingénierie propre (truststore,
modèles hors ligne, backends commutables, stockage embarqué, file sérialisée) plutôt qu'un
contournement fragile. Les résultats RAGAS valident les deux paris techniques centraux,
l'indexation multimodale par description VLM et la recherche hybride RRF, tout en
identifiant avec précision les axes d'amélioration.
